\documentclass[12pt]{article}
\usepackage[utf8]{inputenc}
\usepackage[T1]{fontenc}
\usepackage{hyperref}
\title{Tukumanf\'ed Language Resources for K-12 Curriculum Development}
\author{Compiled by Claude Sonnet 5 (claude-sonnet-5) for Andr\'as Kornai}
\date{}
\begin{document}
\maketitle

\section{Introduction}
Tukumanf\'ed (ISO 639-3 tkf) is described as an extinct Tupian language of Brazil. My search returned no page on Tukumanf\'ed itself, so I cannot confirm its classification, location or documentation.

\section{Primary Educational Resources}
None found.

\section{Linguistic and Reference Materials}
Only general background on the family was found: Wikipedia, ``Tupian languages'': \url{https://en.wikipedia.org/wiki/Tupian_languages}. It says nothing specific to Tukumanf\'ed.

\section{Religious and Community Resources}
None found.

\section{Digital and Audio Resources}
None found.

\section{Curriculum Development Context}
No attestation or descendant community was identified in the search. A heritage curriculum cannot be responsibly built from what I found. Further checks of Glottolog, the ISO 639-3 registry and Brazilian ethnolinguistic sources (searching under alternate spellings and Portuguese names) are the next step.
\end{document}